\chapter{Build and Tooling}
\docsource{docs/design/04-build-tooling.md}

\section{Build Policy}

CloudCropper의 기본 빌드는 CMake 3.24 이상과 vcpkg manifest mode를 기준으로
한다. 단, \texttt{find\_package} first 스타일을 유지해 system package나 Conan
환경에서도 같은 target graph를 해결할 수 있게 둔다.

\begin{longtable}{p{0.32\linewidth}p{0.58\linewidth}}
\toprule
\textbf{항목} & \textbf{결정} \\
\midrule
C++ standard & C++20, compiler extension off. \\
Compilers & GCC 11+, Clang 14+, MSVC VS2022. \\
Dependency manager & vcpkg manifest mode primary. \\
Tests & GoogleTest + CTest. \\
GUI deps & GLFW, OpenGL, Dear ImGui, ImGuizmo, glm. \\
Heavy deps & PCL/VTK/Open3D 계열은 optional/off by default. \\
\bottomrule
\end{longtable}

\section{Target Graph}

\begin{lstlisting}
cloudcropper              executable
  -> cloudcropper::core
  -> cloudcropper::io
  -> cloudcropper::transport
  -> cloudcropper::viewer       only when GUI is ON
\end{lstlisting}

public header는 \ccfile{include/cloudcropper/<module>/}에 두고, private header는
\ccfile{src/<module>/} 옆에 둔다. 각 target은 필요한 dependency만 link하고,
feature flag는 source file과 compile definition을 함께 제어한다.

\section{Feature Flags}

\begin{longtable}{p{0.42\linewidth}p{0.48\linewidth}}
\toprule
\textbf{Option} & \textbf{의미} \\
\midrule
\ccoption{CLOUDCROPPER\_BUILD\_GUI} & viewer target과 GUI app path를 켠다. \\
\ccoption{CLOUDCROPPER\_WITH\_PLY} & PLY reader/writer를 빌드한다. \\
\ccoption{CLOUDCROPPER\_WITH\_PCD} & PCD reader/writer를 빌드한다. \\
\ccoption{CLOUDCROPPER\_WITH\_NPZ} & NPZ/NPY reader/writer를 빌드한다. \\
\ccoption{CLOUDCROPPER\_TRANSPORT\_GZIP} & gzip transport/cache codec을 켠다. \\
\ccoption{CLOUDCROPPER\_BUILD\_TESTS} & GoogleTest/CTest test target을 켠다. \\
\ccoption{CLOUDCROPPER\_ENABLE\_SANITIZERS} & ASan/UBSan interface target을 link한다. \\
\bottomrule
\end{longtable}

\section{Developer Commands}

\begin{lstlisting}[language=bash]
cmake --preset dev
cmake --build --preset dev
ctest --preset dev
\end{lstlisting}

LaTeX 정리본은 애플리케이션 빌드와 독립적이다. 문서 PDF는
\ccfile{docs/organize}에서 다음 명령으로 생성한다.

\begin{lstlisting}[language=bash]
make
\end{lstlisting}

